\section{Effective point output for globally convex polynomials}
\label{sec:global}

The first positive result converts value accuracy into distance accuracy
without a supplied growth modulus. It allows nonunique minima, singular
Hessians, unbounded polyhedra, and variable degree. The polynomial must be
convex on all of Euclidean space.

\begin{theorem}[Global point oracle]\label{thm:global}
Let $f\in\Q[x_1,\ldots,x_n]$ be globally convex and explicitly sparse,
and let $P=\{x:Bx\le b\}$ be a rational polyhedron. Write $I$ for the
binary input length and $D=\max\{2,\deg f\}$. In $\poly(I,D)$ bit work
one can decide emptiness and unboundedness below. Otherwise the minimum is
attained, and one can compute positive rationals $R,\Gamma\ge1$ with
\[
 \log R+\log\Gamma=\poly(I,D)
\]
such that the minimum-norm optimizer $p$ satisfies $\norm p\le R$ and
\begin{equation}\label{eq:global-error}
 \dist(x,S)\le\Gamma\bigl(f(x)-f_*\bigr)^{1/D}
 \quad(x\in P,\ 0\le f(x)-f_*\le1).
\end{equation}
For every $q\ge0$, a feasible rational approximation to this same $p$
within $2^{-q}$ can be computed in $\poly(I,D,q)$ bit work. A simultaneous
objective enclosure of width $2^{-s}$ has work $\poly(I,D,q+s)$.
\end{theorem}

The numerical degree appears in this bound. With unary degree data it is a
polynomial-input-length theorem; with exponentially large binary exponents
it is not. Convexity is a promise, or its certification must be charged.
The empty and zero-variable cases are direct rational decisions.

Small optimizer norms, attainment, unboundedness detection, and value-gap
optimization are already established by \citet{SlotSteurerWiedmer2025}. The additional
output here uses an effective error bound and a quantitative selector
schedule. The radius argument below provides direct rational constants and
preserves sparse evaluation. It also makes explicit which prior structural
facts are needed. The full constants and representation accounting are in
\href{../global-point/theorem.md}{the companion proof}.

\subsection{A rational quadratic lower bound and a radius}

Put $K_D=(D+1)D^{D+2}$. For $a,x\in\R^n$, the Bregman polynomial
\[
 h(t)=f(a+t(x-a))-f(a)-t\nabla f(a)^\top(x-a)
\]
is nonnegative, vanishes at zero, and satisfies $0\le h(t)\le h(1)$ on
$[0,1]$. Lagrange interpolation at $j/D$, $0\le j\le D$, bounds its
second derivative at zero by $K_Dh(1)$. Hence
\begin{equation}\label{eq:tangent-quadratic}
 f(x)\ge f(a)+\nabla f(a)^\top(x-a)
       +K_D^{-1}(x-a)^\top\nabla^2f(a)(x-a).
\end{equation}
Sharper tangent-quadratic bounds are prior; this elementary constant
suffices for the encoding estimates.

Average \eqref{eq:tangent-quadratic} over $a\in[0,1]^n$. The resulting
matrix $M$, vector $\ell$, and scalar $c$ are rational:
\begin{align*}
 M&=\int\nabla^2f(a)\,da,\\
 \ell&=\int\bigl[\nabla f(a)-2K_D^{-1}\nabla^2f(a)a\bigr]\,da,\\
 c&=\int\bigl[f(a)-\nabla f(a)^\top a+
                     K_D^{-1}a^\top\nabla^2f(a)a\bigr]\,da.
\end{align*}
Each integral is evaluated termwise using
$\int a^\alpha da=\prod_i(\alpha_i+1)^{-1}$. Sparse differentiation and
integration have polynomial bit cost in $I,D$. We obtain
\begin{equation}\label{eq:average-lower}
 f(x)\ge c+\ell^\top x+K_D^{-1}x^\top Mx.
\end{equation}

The PSD matrix $M$ has the same kernel as the common kernel of all
Hessians. Indeed, $d^\top Md=0$ forces the continuous nonnegative
polynomial $d^\top\nabla^2f(a)d$ to vanish on the cube, hence everywhere.
PSD then gives $\nabla^2f(a)d=0$. Let $\Pi$ be the rational orthogonal
projection onto $\operatorname{range}M$, and put
$w=-(I_n-\Pi)\nabla f(0)$. Integration along kernel directions gives
\begin{equation}\label{eq:invariance-slope}
 f(x)=f(\Pi x)-w^\top x,\qquad \ell=\Pi\ell-w.
\end{equation}
This nonlinear-subspace and affine-direction distinction is part of the
prior structure theory; no transformed polynomial is expanded here.

Integer-minor bounds give an explicit $\rho>0$ with polynomial bit length
such that $K_D^{-1}x^\top Mx\ge\rho\norm{\Pi x}^2$. For example, clear
denominators with $Q_M$, bound the entries of $Q_MM$ by $C_M\ge1$, and
take
\[
 \rho=\bigl[K_DQ_M(nC_M)^{n-1}\bigr]^{-1}.
\]
The product of the nonzero eigenvalues of an integer PSD matrix is a
positive integer, which proves the claimed bound. If $M=0$, both sides of
the quadratic inequality are zero and the same positive formula is harmless.

Find a rational feasible $a\in P$. The rational linear system
\begin{equation}\label{eq:boundedness-lp}
 w=B^\top\lambda+Mz,\qquad\lambda\ge0
\end{equation}
decides boundedness below. If it is infeasible, Farkas' lemma supplies
$Bd\le0$, $Md=0$, $w^\top d>0$; the feasible ray $a+td$ has objective
decreasing without bound by \eqref{eq:invariance-slope}. If it is feasible,
put $v=Mz$, $A=\norm{\Pi\ell-v}_1$, and $\beta=c-\lambda^\top b$.
Equations \eqref{eq:average-lower}--\eqref{eq:boundedness-lp} give
\begin{equation}\label{eq:quotient-bound}
 f(x)\ge\beta-A\norm{\Pi x}+\rho\norm{\Pi x}^2\qquad(x\in P).
\end{equation}
This is a finite lower bound. On the sublevel $f(x)\le f(a)$, it bounds
$\norm{\Pi x}$ by an explicit polynomial-bit number $T$. It also bounds
$w^\top x$ on both sides: the upper bound follows from
$w^\top x\le\lambda^\top b+v^\top\Pi x$, and the lower bound from
\[
 w^\top x=f(\Pi x)-f(x)
       \ge f(0)-\norm{\nabla f(0)}_1T-f(a).
\]
The one-sided linear constraints alone do not bound $|w^\top x|$.

Stack $M$ and $w^\top$ into a rational matrix $J$. For every point $y$
in this sublevel, the slice
$Z_y=\{x\in P:Jx=Jy\}$ contains an equally good feasible representative
of polynomially bounded logarithmic norm. To see this, clear denominators
in $J$ and in the constraint rows. If their integer entries are bounded by
$C$, the right-hand-side-uniform Hoffman estimate gives
\begin{equation}\label{eq:hoffman-radius}
 \dist(a,Z_y)\le(nC)^{n-1}\norm{J_Z(a-y)},
\end{equation}
where $J_Z$ is the integer-scaled equality matrix. The right-hand side is
bounded using $T$ and the two bounds on $w^\top y$.

For completeness, the Hoffman estimate follows by projecting onto the
slice, retaining at most $n$ independent equality and active inequality
normals, and bounding their least singular value below by
$(nC)^{-(n-1)}$. Their Gram determinant is a positive integer. The signs
of the inequality multipliers permit the equality residual alone to bound
the projection distance for $a\in P$. An irrational equality right-hand
side causes no difficulty.

Equality of $J$-values preserves $f$ by \eqref{eq:invariance-slope}.
A minimizing sequence therefore has equally good representatives in a fixed
closed ball with rational polynomial-bit radius $R$. Compactness proves
attainment. The minimum-norm optimizer has norm at most $R$. This argument
does not presume attainment in order to prove the radius.

\subsection{Sparse coefficient rows control distance}

Write $f(x)=\sum_\alpha f_\alpha x^\alpha$. For every exponent $\gamma$
appearing in a partial derivative, form the row
\[
 N_{\gamma,i}=(\gamma_i+1)f_{\gamma+e_i}.
\]
Missing coefficients are zero. With $s$ original monomials, there are at
most $ns$ rows. The identity
\[
 \nabla f(z)^\top d=\sum_\gamma (Nd)_\gamma z^\gamma
\]
shows that $\ker N$ is the global translation-invariance space. Computing
this coefficient space is prior structure; the following quantitative bound
is what converts it into a point oracle.

Fix the minimum-norm optimizer $p$, any $x\in P$, and write
$d=x-p$, $E=f(x)-f_*\in(0,1]$. The case $E=0$ already has $x\in S$.
The polynomial
\[
 h(u)=f(p+u)-f(p)-\nabla f(p)^\top u
\]
is globally nonnegative. Convexity and constrained optimality give
$0\le h(td)\le E$ for $0\le t\le1$. Univariate interpolation gives
\[
 |h(2td)|\le (D+1)D^D E(1+2|t|)^D.
\]
For $c\in[0,1]^n$, global convexity further gives
\[
 0\le h(c-p+td)\le\tfrac12h(2(c-p))+\tfrac12h(2td).
\]
The first term has a rational coefficient-sum bound depending on $R$ and
$f$. On $0\le t\le E^{-1/D}$, the entire right-hand side has a bound
of polynomial logarithmic size. Differentiating its univariate
interpolation formula at zero yields
\begin{equation}\label{eq:gradient-residual}
 |\nabla f(c)^\top(x-p)|\le C_* E^{1/D}
 \qquad(c\in[0,1]^n),
\end{equation}
with explicitly computable rational $C_*$ and
$\log C_*=\poly(I,D)$. Only $p$ needs a radius bound; $x$ can lie anywhere
in the original polyhedron. If $E=0$, the polynomial along $d$ vanishes
identically and the midpoint bound makes the second polynomial constant,
so \eqref{eq:gradient-residual} holds with zero right-hand side.

Put $r=D-1$ and $B_r=(r+1)2^r r^r$. Tensor Lagrange interpolation shows
that each coefficient of a polynomial of individual degree at most $r$ is
bounded by $B_r^n$ times its supremum on $[0,1]^n$. This conceptual grid
is never enumerated: its only algorithmic consequence is the integer
$B_r^n$, whose bit length is $O(nD\log(D+1))$. Consequently
\[
 |(N(x-p))_\gamma|\le B_r^n C_* E^{1/D}.
\]
At zero gap this proves that every optimizer has the same $N$-value.
Conversely, invariance preserves the objective, giving
\begin{equation}\label{eq:optimizer-slice}
 S=\{x\in P:Nx=Np\}.
\end{equation}

Clear denominators with $Q_N$, and bound the integer entries of $Q_NN$
and the scaled constraint rows by $C_N\ge1$. The Hoffman estimate gives
\eqref{eq:global-error}, for example with
\[
 \Gamma=\max\{1,(nC_N)^{n-1}Q_N\max(1,\#\text{rows of }N)
                                      B_r^n C_*\}.
\]
Every operation uses sparse input and rational numbers of polynomial bit
length. A constant objective gives an empty coefficient matrix and $S=P$,
which is covered without a rank assumption.

\subsection{Selecting one fixed optimizer}

Let $\varepsilon=2^{-q}$ and set
\[
 \tau=\frac{\varepsilon^{2D-2}}
              {4\,8^{D-1}R^D\Gamma^D},\qquad
 \eta=\frac{\tau\varepsilon^2}{4}.
\]
The minimizer $x_\tau$ of $f(x)+\tau\norm x^2$ over $P$ exists,
is unique, and satisfies $\norm{x_\tau}\le\norm p\le R$ by comparison
with $p$. Let $s$ be a nearest optimizer to $x_\tau$, and put
$e=\norm{s-x_\tau}$. Since $e\le2R$, comparison with $s$ and
\eqref{eq:global-error} give
\[
 e^D/\Gamma^D\le f(x_\tau)-f_*
          \le\tau e(2\norm{x_\tau}+e)\le4R\tau e.
\]
Thus $e\le\varepsilon^2/(8R)$. The projection inequality
$p^\top(s-p)\ge0$ and $\norm{x_\tau}\le\norm p$ imply
\[
 \norm{x_\tau-p}^2\le2R e\le\varepsilon^2/4.
\]
The coefficient $4R$ allows the nearest optimizer $s$ to lie outside the
radius box.

Optimize the regularized polynomial over
$Q=P\cap[-R,R]^n$ to certified gap $\eta$. Its modulus is $2\tau$,
so the returned feasible rational point is within $\varepsilon/2$ of
$x_\tau$, hence within $\varepsilon$ of $p$. The required coefficient
and accuracy lengths are polynomial in $I,D,q$. The convex value routine
retains the original sparse polynomial during affine-hull reduction and
evaluates the composition by the chain rule, rather than expanding it.
Classical rational feasibility repair and value certification then apply.

This completes the proof of Theorem~\ref{thm:global}. The theorem does not
decide exact active constraints or exact coordinate signs. Those require a
different output guarantee even though arbitrarily accurate rational points
are available.
